\section{Accountability: who is responsible when it is wrong?}
\label{sec:accountability}

A clinician must know where responsibility sits. The seventh question asks who
answers when something goes wrong, what is documented and by whom, and how
liability is allocated across the human and the system.

[DRAFTING INSTRUCTIONS: in the full-framework, describe explicit roles, the signed
standard operating procedure, and the documentation chain, grounded in the
author's documentation and regulatory-adaption lineage in
\texttt{review/references/author\_works.bib}
(\texttt{Kawchak2026PhysicalAIOncologyClinical},
\texttt{Kawchak2026Adaption21CFRPart},
\texttt{Kawchak2026AdaptionICHHarmonisedGuideline}). Reproduce Mermaid
\texttt{adoption/mermaid/diagrams/13-accountability-raci.md} as colored cfwfig
TikZ. Carry the KEY TERMS ``standard operating procedure'' and ``audit trail''.
Build a body-width RACI table (clinician, sponsor, developer, IRB, auditor) by
function.]

\begin{cfwfig}
\node[cfact] (a) at (0,0) {Clinician}; \node[cftwo] (b) at (0,-1.4) {Sponsor};
\node[cftwo] (c) at (0,-2.8) {Developer}; \node[cfhope] (d) at (5.0,-1.4) {Signed SOP};
\node[cfone] (e) at (9.4,-1.4) {Auditor};
\draw[cfedge] (a)--(d); \draw[cfedge] (b)--(d); \draw[cfedge] (c)--(d); \draw[cfedge] (d)--(e);
\end{cfwfig}
\vspace{-0.2cm}
\cfwcaption{Figure 9. Where accountability sits (placeholder; expanded from Mermaid
13 in the full-framework).}

Accountability is answered by drawing it: named roles, a signed procedure, and a
documentation chain that an auditor can follow from action to owner.
